\documentclass[11pt]{article}
\usepackage[a4paper,margin=22mm]{geometry}
\usepackage{amsmath,amssymb,amsthm,mathtools,bm}
\usepackage{slashed}
\usepackage{mathrsfs}
\usepackage{booktabs,longtable,array}
\usepackage{enumitem}
\usepackage{xcolor}
\usepackage{microtype}
\usepackage[hidelinks]{hyperref}

\setlength{\emergencystretch}{3em}
\newtheorem{theorem}{Theorem}[section]
\newtheorem{proposition}[theorem]{Proposition}
\newtheorem{lemma}[theorem]{Lemma}
\newtheorem{corollary}[theorem]{Corollary}
\theoremstyle{definition}
\newtheorem{definition}[theorem]{Definition}
\newtheorem{gate}[theorem]{Fail-closed gate}
\theoremstyle{remark}
\newtheorem{remark}[theorem]{Remark}

\newcommand{\dd}{\mathrm d}
\newcommand{\ii}{\mathrm i}
\newcommand{\ZZ}{\mathbb Z}
\newcommand{\RR}{\mathbb R}
\newcommand{\CC}{\mathbb C}
\newcommand{\cH}{\mathcal H}
\newcommand{\cD}{\mathcal D}
\newcommand{\cR}{\mathcal R}
\newcommand{\cF}{\mathcal F}
\newcommand{\Tr}{\operatorname{Tr}}
\newcommand{\tr}{\operatorname{tr}}
\newcommand{\ind}{\operatorname{ind}}
\newcommand{\SF}{\operatorname{SF}}
\newcommand{\FP}{\operatorname{FP}}
\newcommand{\Arf}{\operatorname{Arf}}
\newcommand{\status}[1]{\textcolor{blue!60!black}{\textsf{#1}}}
\newcolumntype{L}[1]{>{\raggedright\arraybackslash}p{#1}}

\title{AP-E7: A Common Quadratic (Sp(4)) APS/PV Regulator\\
Exact Pure Heavy-Threshold Phase, Two Target Tests,\\
and a Determinant-Line No-Go}
\author{Route F independent research note}
\date{17 July 2026}

\begin{document}
\maketitle

\begin{abstract}
This note repairs a precise omission of AP-E6.  For the candidate simply
connected group, denoted \(Sp(4)\) in the physics convention and
\(Sp(2)\) in the mathematics convention, we specify one common
background-field quadratic regulator for fermions, gauge fields,
    the adjoint and fundamental scalars, and Faddeev--Popov ghosts.  The
    Hilbert spaces, Sobolev domains, Euclidean Clifford and reference real structures,
gauge complex, Pauli--Villars superdeterminants, finite local subtraction,
five-dimensional mass interpolation, spectral cut, and APS domains are all
explicit.  This closes
\texttt{common\_quadratic\_aps\_pv\_regulator\_specified}.  It does not
construct a nonperturbative continuum gauge measure: Gribov/global gauge
slicing, interacting removal of the cutoff, and reflection positivity remain
open.

As a constructive negative control, we also give a compact-link overlap
framework and a continuum Higgs-orbit source that are mutually
Ginsparg--Wilson compatible.  We do \emph{not} write the finite-lattice
Dirac--Yukawa operator or its five-dimensional target kernel, so this is not a
completed overlap lift and no overlap eta pair is claimed as computed.  A
copy-counting diagnostic predicts the conditional pair \((+1,+1)\) if such a
lift, the standard determinant orientation, and the no-stacked-torsion
convention are supplied.  This is new conditional Higgsed-sector data, not a
consequence of the original action.  A Schur-lemma theorem proves only that
an exact separated rank-two light projector cannot extend continuously and
\(Sp(4)\)-equivariantly to the restored point \(\Sigma=0\).

The same regulator gives a new exact threshold result.  For one flavour the
negative-mass heavy bundle after
\(Sp(4)\to SU(2)_c\times U(1)_H\) is
\(E_h=(V_c\otimes L^{-1})\oplus\bm1\), and
\[
 I_h=\ind D^+_{E_h}
 =\int_{M_4}\left[-c_2(V_c)+x^2-\frac18p_1(TM_4)\right],
 \qquad x=c_1(L).
\]
The two Dirac flavours therefore contribute
\(\exp(2\pi\ii I_h)=1\) on every closed spin four-manifold.  Thus the
constant-mass pure unbroken-gauge plus gravitational APS threshold is
computed and is removable by the displayed integral local counterterm.  On
the two target representatives
\(G_3=S^1_{\rm R}\times S^3\) and
\(G_2=T^2_{\rm RR}\times S^2\), the same pure component is explicitly
\((+1,+1)\).

This does not determine the full target phases.  The microscopic action has
no map from \((g,s)\in S^3\times S^2\) to a full all-branch chiral
Dirac--Yukawa/PV family: \(g\) is composite and no \(s\)-dependent fermion
mass is declared.  Consequently there is no pullback determinant line over
\(\operatorname{Map}(M_4,S^3\times S^2)\).  We prove non-identifiability by
stacking the two independent zero-curvature spin characters
\((-1)^{\nu_3}\) and \((-1)^{\nu_2}\), which leaves every local symbol and
the pure threshold unchanged while flipping either target sign.  Moreover,
two controlled ansaetze have been evaluated: a uniform exact flavour mass map
on the complete \(\bm5\) has mixed coefficient \(+2-2=0\), whereas a
light-only emergent flavour assignment has no declared deep-UV target family.
These cases do not exhaust continuous equivariant non-projector couplings.
Hence the present work does not give an all-scale \(k=+2\) matching theorem.
No portal is authorized.
\end{abstract}

\tableofcontents

\section{Claim boundary and field content}

The candidate field list and branch convention are
\begin{equation}
 \Psi_i\in\bm5\ \hbox{Dirac},\quad i=1,2;\qquad
 \Phi_I\in\bm4\ \hbox{complex},\quad I=1,2;\qquad
 \Sigma\in\bm{10}\ \hbox{real}.
 \label{eq:field-list}
\end{equation}
With the Hermitian generator \(H=-2\ii T_H^3\),
\begin{equation}
 \bm5\longrightarrow\bm2_{+1}\oplus\bm2_{-1}\oplus\bm1_0,
 \qquad
 \bm4\longrightarrow\bm2_0\oplus\bm1_{+1}\oplus\bm1_{-1}.
 \label{eq:branching}
\end{equation}
The unbroken group is the direct product
\(SU(2)_c\times U(1)_H\), not its diagonal quotient.  These facts use the
injective \(Spin(4)\subset Spin(5)\) embedding and the standard low
representations \cite{Slansky1981}.

Four statements are separated throughout.
\begin{enumerate}[label=(\roman*)]
\item A common \emph{quadratic} APS/PV regulator can be an exact elliptic
  construction around a nondegenerate smooth background.
\item Such a construction is not a rigorous nonperturbative measure for an
  interacting four-dimensional gauge theory.
\item A pure unbroken-gauge/gravitational mass-sign phase is a determinant
  ratio on bundles over spacetime.
\item A target torsion phase additionally requires a map from target fields
  to a regulated Fredholm family and its determinant line.
\end{enumerate}
The last distinction is the determinant-line version of the APS and
Dai--Freed framework \cite{AtiyahPatodiSinger1975,BismutFreed1986,DaiFreed1994}.
Recent bordism analyses of \(Sp\) QCD and mod-two \(Sp(n)\) anomalies support
the same separation between local polynomials and torsion characters
\cite{Saito2024,SaitoTachikawa2025}.

\section{Euclidean Hilbert spaces, gamma matrices, and reality}

\subsection{Geometric data and domains}

Let \(M_4\) be a closed connected Riemannian spin four-manifold, let
\(S=S^+\oplus S^-\) be its complex spinor bundle, and let
\(P\to M_4\) be a principal \(Sp(2)\)-bundle.  Write \(E_R=P\times_R R\).
All connections are smooth and anti-Hermitian.  The basic Hilbert spaces are
\begin{align}
 \cH_F&=L^2(M_4,S\otimes E_{\bm5}\otimes\CC^2_f),
 &\operatorname{Dom}\slashed D_A&=H^1(M_4,S\otimes E_{\bm5}\otimes\CC^2_f),
 \label{eq:fermion-domain}\\
 \cH_B&=L^2(M_4,\mathcal B),
 &\operatorname{Dom}\Delta_B&=H^2(M_4,\mathcal B),
 \label{eq:boson-domain}\\
 \cH_{\rm gh}&=L^2(M_4,\operatorname{ad}P\otimes\CC),
 &\operatorname{Dom}\Delta_{\rm gh}&=H^2(M_4,\operatorname{ad}P\otimes\CC),
 \label{eq:ghost-domain}
\end{align}
where
\begin{equation}
 \mathcal B=
 \bigl(T^*M_4\otimes\operatorname{ad}P\bigr)
 \oplus\operatorname{ad}P
 \oplus\bigl(E_{\bm4}^{\oplus2}\bigr)_{\RR}.
 \label{eq:boson-bundle}
\end{equation}
The last summand is realified because the scalar Gaussian is a real bosonic
Gaussian.  Compactness and ellipticity give discrete spectra with finite
multiplicity.  A prime on a determinant will remove the finite-dimensional
kernel and place its orientation in the determinant line.

\subsection{Clifford and real structures}

Choose Hermitian Euclidean gamma matrices
\begin{equation}
 \gamma_1=\sigma_1\otimes\sigma_1,\quad
 \gamma_2=\sigma_1\otimes\sigma_2,\quad
 \gamma_3=\sigma_1\otimes\sigma_3,\quad
 \gamma_4=\sigma_2\otimes\bm1,
 \qquad
 \gamma_5=\gamma_1\gamma_2\gamma_3\gamma_4.
 \label{eq:gamma-choice}
\end{equation}
Thus
\begin{equation}
 \gamma_\mu^\dagger=\gamma_\mu,\qquad
 \{\gamma_\mu,\gamma_\nu\}=2\delta_{\mu\nu},\qquad
 \gamma_5^\dagger=\gamma_5,\qquad \gamma_5^2=1.
 \label{eq:clifford-checks}
\end{equation}
Put \(B=\gamma_2\gamma_4\) and let \(K\) denote componentwise complex
conjugation.  Direct multiplication gives
\begin{equation}
 B\gamma_\mu^*B^{-1}=\gamma_\mu,\qquad BB^*=-\bm1.
 \label{eq:spin-reality}
\end{equation}
The \(\bm5\) is the real vector representation of \(Spin(5)\), so its
bundle has a parallel conjugation \(K_{\bm5}^2=+1\).  Therefore
\begin{equation}
 J=(BK)\otimes K_{\bm5},\qquad J^2=-1
 \label{eq:total-reality}
\end{equation}
is the quaternionic real structure on the one-flavour spinor bundle.  It
commutes with the covariant Dirac operator and a real scalar reference mass.
It does \emph{not} commute with a fixed nonzero adjoint Higgs mass in the
convention below.  Indeed, in a real vector basis the anti-Hermitian
generators \(T^a_{\bm5}\) are real antisymmetric, and hence
\begin{equation}
 K_{\bm5}\widehat\Sigma_{\bm5}K_{\bm5}^{-1}
 =-\widehat\Sigma_{\bm5},\qquad
 J\mathcal K_{A,\Sigma}J^{-1}=\mathcal K_{A,-\Sigma}.
 \label{eq:adjoint-reality-odd}
\end{equation}
Thus \(J^2=-1\) gives Kramers pairing only in the scalar-reference sector
(or for a family involution that also sends \(\Sigma\mapsto-\Sigma\)), not at
a fixed generic Higgs background.  Positivity below instead uses the
elementary two-flavour determinant square and gamma-five Hermiticity; no
fixed-adjoint-mass Kramers claim is made.

For anti-Hermitian \(A\), \(\slashed D_A^\dagger=-\slashed D_A\).  Define
\begin{equation}
 \widehat\Sigma_{\bm5}=-\ii\Sigma^aT^a_{\bm5},\qquad
 \mathcal K_{A,\Sigma}=\slashed D_A+m+y\widehat\Sigma_{\bm5},\qquad
 H_{A,\Sigma}=\gamma_5\mathcal K_{A,\Sigma}.
 \label{eq:fermion-H}
\end{equation}
Then
\begin{equation}
 \mathcal K^\dagger=\gamma_5\mathcal K\gamma_5,\qquad
 H^\dagger=H,\qquad \operatorname{Dom}H=H^1.
 \label{eq:gamma5-hermiticity}
\end{equation}
The light mass is kept at \(M_+=\mu>0\) during every determinant
calculation; the limit \(\mu\downarrow0\) is taken only after zero modes have
been separated.

\section{The full background-field quadratic complex}

\subsection{Gauge action and gauge fixing}

Fix a smooth background \((\bar A,\bar\Sigma,\bar\Phi)\) and write the
fluctuation as \(b=(a,\sigma,\varphi)\in\Gamma(\mathcal B)\).  At vanishing
background fermions the infinitesimal gauge-action operator is
\begin{equation}
 \cR_{\bar X}\epsilon=
 \bigl(D_{\bar A}\epsilon,
       [\bar\Sigma,\epsilon],
       -\epsilon\bar\Phi\bigr),
 \qquad
 \cR_{\bar X}:H^1(\operatorname{ad}P)\longrightarrow L^2(\mathcal B).
 \label{eq:gauge-action}
\end{equation}
The notation \(-\epsilon\bar\Phi\) includes both fundamental copies and
realification.  The background Feynman gauge functional is the exact adjoint
\begin{equation}
 \cF_{\bar X}(b)=\cR_{\bar X}^{\dagger}b,
 \qquad
 S_{\rm gf}^{(2)}=\frac12\|\cF_{\bar X}(b)\|^2.
 \label{eq:gauge-fixing}
\end{equation}
Let \(S_{\rm bos}\) be the AP-E6 Yang--Mills plus covariant kinetic and
scalar-potential action.  Its gauge-fixed Hessian and ghost operator are
defined without dropping mixing terms by
\begin{align}
 \Delta_B&=S_{\rm bos}^{(2)}[\bar X]+\cR_{\bar X}\cR_{\bar X}^{\dagger},
 \label{eq:exact-boson-hessian}\\
 \Delta_{\rm gh}&=\cR_{\bar X}^{\dagger}\cR_{\bar X}.
 \label{eq:exact-ghost-hessian}
\end{align}
This definition contains the gauge--adjoint--fundamental mixing and the full
matrix of second derivatives of the declared potential.  Their principal
symbols are
\begin{equation}
 \sigma_2(\Delta_B)(x,\xi)=|\xi|^2\bm1_{\mathcal B},
 \qquad
 \sigma_2(\Delta_{\rm gh})(x,\xi)=|\xi|^2\bm1_{\operatorname{ad}P}.
 \label{eq:principal-symbols}
\end{equation}
Hence both are Laplace type and essentially self-adjoint on their stated
Sobolev domains.  Background stabilizers are projected from both operators.
If \(\Delta_B\) has negative modes, the Gaussian contour is the steepest
descent contour fixed by the upper lateral Agmon cut
\(\theta_B=\pi-\varepsilon\), \(0<\varepsilon\ll1\); such a saddle is not
called stable.  This convention defines the quadratic determinant but does
not turn a negative mode into a physical vacuum.

\subsection{What is and is not included}

Equations~\eqref{eq:exact-boson-hessian}--\eqref{eq:exact-ghost-hessian}
are preferable to listing only diagonal blocks: the latter would silently
drop the Higgs-mechanism mixing.  The quadratic functional is
\begin{equation}
 S^{(2)}=\frac12\langle b,\Delta_Bb\rangle
 +\langle\bar c,\Delta_{\rm gh}c\rangle
 +\sum_{i=1}^2\langle\bar\Psi_i,\mathcal K\Psi_i\rangle.
 \label{eq:full-quadratic-action}
\end{equation}
This is the entire gauge--scalar--ghost--fermion quadratic complex about the
declared background.  Fermion-loop contributions to an effective bosonic
Hessian arise after integrating Eq.~\eqref{eq:full-quadratic-action}; they
are not part of the bare Gaussian operator and are not claimed here.

\section{One common Pauli--Villars prescription}

\subsection{Superdeterminant and regulator statistics}

Use the four-node finite difference
\begin{equation}
 c_a=(1,-3,3,-1),\qquad \Lambda_a^2=a\Lambda^2,\qquad a=0,1,2,3.
 \label{eq:pv-data}
\end{equation}
It obeys
\begin{equation}
 \sum_ac_a a^r=0\quad(r=0,1,2),
 \qquad
 \sum_ac_aa^3=-6.
 \label{eq:pv-moments}
\end{equation}
For any positive Laplace-type eigenvalue \(\lambda\),
\begin{equation}
 \sum_ac_a\log(\lambda+a\Lambda^2)
 =-\frac{2\Lambda^6}{\lambda^3}+O(\lambda^{-4}),
 \label{eq:pv-asymptotic}
\end{equation}
so the first three ultraviolet moments cancel.  Define
\begin{align}
 Z_{\rm quad}^{\rm PV}(\Lambda)
 :={}&\tau_F
 \prod_{a=0}^3
 \det{}'_{\theta_B}(\Delta_B+\Lambda_a^2)^{-c_a/2}
 \det{}'(\Delta_{\rm gh}+\Lambda_a^2)^{c_a}
 \notag\\
 &\times
 \det{}'(H^2+\Lambda_a^2)^{N_f c_a/2},
 \qquad N_f=2.
 \label{eq:common-pv}
\end{align}
Here \(\tau_F\) is the fermion determinant-line phase defined in the next
section.  The last determinant is only its positive amplitude.  It can be
implemented with an auxiliary Pauli doublet
\begin{equation}
 \widehat H_a=H\otimes\sigma_3+\Lambda_a\bm1\otimes\sigma_1,
 \qquad \widehat H_a^2=H^2+\Lambda_a^2,
 \label{eq:pv-doublet}
\end{equation}
whose spectrum is phase-paired.  Thus amplitude subtraction cannot
silently choose a torsion phase.

Equation~\eqref{eq:common-pv} is first of all a zeta-determinant
prescription.  Integer powers in complex sectors admit the usual commuting
or Berezin auxiliary-field representation, with multiplicity \(|c_a|\).
For the real bosonic sector the half-integer power is instead the spectral-cut
square root of the determinant.  We do not claim that a positive half-integer
power of a symmetric second-order operator is a literal real-Grassmann
Gaussian: such a local realization would require an additional
antisymmetric first-order Pfaffian construction, which is not supplied here.
The three elliptic determinant families are:
\begin{center}
\begin{tabular}{L{0.24\textwidth}L{0.27\textwidth}L{0.36\textwidth}}
\toprule
sector & operator at node \(a\) & determinant exponent\\
\midrule
real gauge plus scalars & \(\Delta_B+a\Lambda^2\) & \(-c_a/2\)\\
complex FP ghost & \(\Delta_{\rm gh}+a\Lambda^2\) & \(+c_a\)\\
two Dirac amplitudes & \(H^2+a\Lambda^2\) & \(N_fc_a/2=c_a\)\\
\bottomrule
\end{tabular}
\end{center}
Every mass is a scalar endomorphism, so the regulator preserves background
\(Sp(4)\) covariance.  BRST covariance at quadratic order follows because
the boson and ghost operators are built from the same \(\cR_{\bar X}\).

\subsection{Finite local subtraction}

The moment equations remove power terms but leave local logarithmic terms.
Let \(A_4(\bar X)\) be the total integrated Seeley--DeWitt coefficient of
the supercomplex in Eq.~\eqref{eq:common-pv}.  At a fixed scale \(\mu_R\)
we define the finite scheme
\begin{equation}
 \log Z_{\rm quad}^{\rm ren}
 =\FP_{\Lambda\to\infty}
 \left[\log Z_{\rm quad}^{\rm PV}(\Lambda)
       -A_4(\bar X)\log\frac{\Lambda^2}{\mu_R^2}
       -S_{\rm pow}^{\rm loc}(\Lambda;\bar X)\right].
 \label{eq:finite-part}
\end{equation}
The finite local convention is fixed by
\begin{equation}
 \theta_c=\theta_H=\theta_{\rm grav}=0
 \quad\hbox{for positive reference masses},
 \qquad
 Z_{\rm quad}^{\rm ren}[0,VH,0]>0,
 \label{eq:finite-convention}
\end{equation}
and by adding no separate target torsion invertible theory.  This fixes all
local de Rham counterterms of the quadratic reference problem.  It cannot
fix a pullback determinant line over a target space for which no pullback
map has been specified.

\begin{proposition}[Quadratic regulator closure]
On a smooth compact background with projected stabilizer and a declared
Agmon contour for every negative bosonic mode,
Eqs.~\eqref{eq:fermion-domain}--\eqref{eq:finite-convention} define the full
common quadratic fermion--PV--gauge--scalar--ghost amplitude.  Together with
the phase construction below,
\begin{equation}
 \texttt{common\_quadratic\_aps\_pv\_regulator\_specified=true}.
 \label{eq:quadratic-regulator-true}
\end{equation}
This statement is deliberately weaker than a nonperturbative continuum
measure.
\end{proposition}

\section{Five-dimensional interpolation, spectral cut, and APS domain}

\subsection{The common cylinder operator}

Let \(Y_5=M_4\times[0,1]\) with product geometry near both ends.  Choose a
smooth path of backgrounds \(X(t)=(A(t),\Sigma(t),\Phi(t))\), constant near
\(t=0,1\), and a Hermitian fermion mass endomorphism \(\mathcal M(t)\).
The positive PV reference is
\begin{equation}
 H_0=\gamma_5(\slashed D_{A(0)}+M_*),\qquad M_*>0,
 \label{eq:positive-reference}
\end{equation}
and the endpoint is
\begin{equation}
 H_1=\gamma_5(\slashed D_{A(1)}+m+y\widehat\Sigma(1)).
 \label{eq:physical-endpoint}
\end{equation}
The simplest mass interpolation is affine away from collars,
\begin{equation}
 \mathcal M(t)=(1-f(t))M_*+f(t)(m+y\widehat\Sigma(t)),
 \quad f=0\ \hbox{near }0,\quad f=1\ \hbox{near }1.
 \label{eq:mass-interpolation}
\end{equation}
No gap is assumed along the path; crossings are the information measured by
spectral flow.

For the physical determinant-sign ratio we require gapped endpoints,
\begin{equation}
 0\notin\operatorname{spec}(H_0)\cup\operatorname{spec}(H_1),
 \label{eq:zero-cut-gap}
\end{equation}
and fix the physical spectral cut to \(\alpha=0\).  Crossings through zero in
the open cylinder remain allowed.  For the separate mathematical statement
about a Fredholm index, one may instead choose any real cut in the common
endpoint resolvent,
\(\alpha\notin\operatorname{spec}(H_0)\cup\operatorname{spec}(H_1)\), and
write \(P_I(H)\) for the spectral projector of \(H\) on \(I\).  The Fredholm
operator
\begin{equation}
 \cD_5^+=\partial_t+H(t)
 \label{eq:five-d-operator}
\end{equation}
has the explicit domain
\begin{equation}
 \begin{split}
 \operatorname{Dom}_{\rm APS,\alpha}(\cD_5^+)
 =\{u\in H^1(Y_5):\ &P_{[\alpha,\infty)}(H_0)u|_{t=0}=0,\\
 &P_{(-\infty,\alpha)}(H_1)u|_{t=1}=0\}.
 \end{split}
 \label{eq:aps-domain}
\end{equation}
Its adjoint is \(-\partial_t+H(t)\) with complementary boundary projectors.
The self-adjoint double is
\begin{equation}
 \mathbb D_{5,\alpha}=
 \begin{pmatrix}0&(\cD_5^+)^*\\ \cD_5^+&0\end{pmatrix}
 \label{eq:selfadjoint-double}
\end{equation}
on the direct sum of these two APS domains.  This gives a literal Hilbert
space and domain, rather than a schematic ``five-dimensional eta
operator.''

\subsection{Reduced eta and determinant phase}

For a self-adjoint elliptic endpoint operator define
\begin{equation}
 \eta_\alpha(H;s)=\sum_{\lambda\ne\alpha}
 \operatorname{sgn}(\lambda-\alpha)|\lambda-\alpha|^{-s},
 \qquad
 \bar\eta_\alpha(H)=\frac{\eta_\alpha(H;0)+\dim\ker(H-\alpha)}2.
 \label{eq:reduced-eta}
\end{equation}
For every such auxiliary cut, the APS theorem identifies
\begin{equation}
 \ind_{\rm APS,\alpha}(\cD_5^+)=\SF_\alpha\{H(t)\}.
 \label{eq:sf-index}
\end{equation}
For a real mass interpolation with the gapped endpoints
\eqref{eq:zero-cut-gap}, the regulated physical determinant-sign ratio is
\begin{equation}
 \tau_F[X(t)]
 =\exp\!\left(\ii\pi N_f\SF_0\{H(t)\}\right).
 \label{eq:fermion-phase}
\end{equation}
The arbitrary \(\alpha\) in Eq.~\eqref{eq:sf-index} is therefore an index
bookkeeping device and must not be substituted into
Eq.~\eqref{eq:fermion-phase}.  A formulation that moves the physical cut
would require the corresponding endpoint spectral-section orientation (or
reduced-eta correction); that extra datum is not used here.  Equivalent
zero-cut reduced-eta formulae differ by the local index density already fixed
in Eq.~\eqref{eq:finite-convention}.  APS gluing and the Bismut--Freed
holonomy theorem make the resulting zero-cut section compatible with
concatenating paths
\cite{AtiyahPatodiSinger1975,BismutFreed1986,DaiFreed1994}.

For a general complex chiral mass family the phase is a section of the
pulled-back determinant line, not simply a sign.  The Dai--Freed prescription
still applies, but only after that family and its endpoint bundles have been
defined \cite{Witten2016,WittenYonekura2019}.

\section{A conditional GW-compatible Higgsed source}

\subsection{Finite compact-link regulator}

There is a useful regulator framework that goes beyond the continuum Gaussian
without pretending to solve the continuum limit or already supplying the
target Yukawa operator.  On a finite four-dimensional
hypercubic lattice take compact \(Sp(4)\) links \(U_\ell\) with Haar measure,
the Wilson plaquette action, and site fields \(\Sigma_x,\Phi_x\) with the
gauge-covariant nearest-neighbour discretization of the declared scalar
action.  The finite bosonic integral is gauge invariant and requires no
gauge fixing or Faddeev--Popov ghosts.  Its perturbative expansion is
matched to the continuum background complex of Section~3.

Let
\begin{equation}
 H_W(U)=\gamma_5(D_W(U)-m_0),\qquad 0<m_0<2,
 \qquad
 D_{\rm ov}=\frac1a\left(1+\gamma_5\operatorname{sgn}H_W\right).
 \label{eq:overlap}
\end{equation}
At finite \(a\), chiral sources use the Ginsparg--Wilson involution and
projectors
\begin{equation}
 \widehat\gamma_5=\gamma_5(1-aD_{\rm ov}),\qquad
 \widehat P_{L,R}=\frac{1\mp\widehat\gamma_5}{2};
 \label{eq:gw-projectors}
\end{equation}
the un-hatted \(P_{L,R}\) in the continuum mass formula below denote their
\(a\to0\) limits.  This keeps the finite-lattice source compatible with the
overlap chiral decomposition.
On the admissible-link locus, \(H_W\) is gapped and the sign function is
local.  A five-dimensional domain-wall determinant divided by its
Pauli--Villars field of dimensionless mass \(m_{\rm PV}=1\) converges as
\(L_s\to\infty\) to the overlap determinant.  Domain-wall eta invariants
realize the APS and mod-two APS indices, while the massive Wilson eta
invariant agrees with the continuum index at sufficiently small lattice
spacing \cite{FukayaOnogiYamaguchi2017,FukayaEtAl2022,AokiEtAl2025}.
These facts provide a finite regulator for a Dirac operator once that operator
and its domain-wall interpolation have been specified.  They do not by
themselves define the target-dependent Dirac--Yukawa family used below.

\subsection{Spectral projectors and an explicit target mass}

Restrict now to the exact Higgs orbit on which
\begin{equation}
 h=\frac{\widehat\Sigma_{\bm5}}V,
 \qquad \operatorname{spec}h=(+1,+1,-1,-1,0),
 \qquad h^3=h.
 \label{eq:higgs-h}
\end{equation}
The covariant orthogonal projectors are
\begin{equation}
 P_+=\frac{h^2+h}{2},\qquad
 P_-=\frac{h^2-h}{2},\qquad
 P_0=1-h^2.
 \label{eq:higgs-projectors}
\end{equation}
They obey
\begin{equation}
 P_rP_s=\delta_{rs}P_r,\qquad
 P_++P_-+P_0=1,\qquad
 (\operatorname{rank}P_+,\operatorname{rank}P_-,\operatorname{rank}P_0)
 =(2,2,1).
 \label{eq:projector-checks}
\end{equation}
Away from the exact orbit but before a spectral gap closes, the same
projectors are defined by Riesz contours around the three separated spectral
clusters.  They cease to exist as branch projectors at the restored locus.

Let \(g_x\in SU(2)_f\) be a declared Hubbard--Stratonovich chiral mass
source.  Let \(A_H(s)\) be the pullback Hopf connection for
\(s:M_4\to S^2\), embedded into the links through
\(\exp(\ii A_H H)\).  The continuum Higgsed all-branch source is
\begin{equation}
 \begin{split}
 \mathcal M_{\rm H}(g,s;\Sigma)
 ={}&P_+\otimes\mu\bigl(P_Lg+P_Rg^\dagger\bigr)
 -2yV P_-\otimes\bm1_f-yV P_0\otimes\bm1_f,\\
 &P_L=\frac{1-\gamma_5}{2},\qquad P_R=\frac{1+\gamma_5}{2}.
 \end{split}
 \label{eq:higgsed-target-mass}
\end{equation}
Here \(s\) is intended to enter the covariant Wilson/overlap operator through
its Hopf link, not as an undeclared direct Yukawa.  Equations
\eqref{eq:overlap}--\eqref{eq:higgsed-target-mass} specify a
GW-compatible continuum-limit source on the gapped Higgsed locus.  They do
not specify the finite operator: the ordinary projectors acting on
\(\bar\Psi\), the hatted projectors acting on \(\Psi\), determinant
orientation, and the corresponding five-dimensional target kernel remain to
be written.  Consequently this note does not claim a completed overlap
determinant line or a computed overlap eta invariant.

There is nevertheless a useful conditional copy-counting diagnostic.  On
\(G_3\), the unit \(S^3\) mass
winding is repeated \(\operatorname{rank}P_+=2\) times, hence its mod-two
index is zero.  On \(G_2\), the light charged sector has
\(2\cdot2\cdot|c_1|=4\) Hopf-flux zero-mode copies; the charged heavy sector
adds another four, and the neutral sector adds none.  Therefore, assuming a
finite GW implementation with the displayed continuum limit, standard
copy-additivity and orientation, positive \(m_{\rm PV}=1\), and no stacked
\(I_3,I_2\) sector, the predicted pair is
\begin{equation}
 \left(\tau_{\rm H,ov}(G_3),\tau_{\rm H,ov}(G_2)\right)=(+1,+1).
 \label{eq:higgsed-overlap-pair}
\end{equation}
This is a conditional parity prediction, not an eta computation or a unique
phase extracted from the original deep-UV action.

\begin{gate}[Higgsed construction versus UV claim]
The exact distinction is
\begin{equation}
 \begin{gathered}
 \texttt{conditional\_higgsed\_gw\_source\_specified=true},\\
 \texttt{finite\_higgsed\_overlap\_dirac\_yukawa\_operator\_defined=false},\\
 \texttt{five\_dimensional\_higgsed\_target\_kernel\_defined=false},\\
 \texttt{minimal\_higgsed\_overlap\_eta\_pair\_computed=false},\\
 \texttt{conditional\_higgsed\_parity\_predicts\_plus\_plus=true},\\
 \texttt{original\_deep\_sp4\_action\_uniquely\_selects\_target\_pair=false}.
 \end{gathered}
 \label{eq:higgsed-lift-gates}
\end{equation}
\end{gate}

\section{The two target generators: what the common regulator computes}

\subsection{Bordism representatives and abstract characters}

Let \(X=S^3\times S^2\).  The two reduced spin-bordism generators are
\begin{align}
 G_3&=(S^1_{\rm R}\times S^3,\operatorname{pr}_{S^3},{\rm pt}),
 \label{eq:G3}\\
 G_2&=(T^2_{\rm RR}\times S^2,{\rm pt},\operatorname{pr}_{S^2}).
 \label{eq:G2}
\end{align}
For regular values define
\begin{equation}
 \nu_3(M,g)=\dim_{\RR}\ker D_{g^{-1}(p_3)}\pmod2,
 \qquad
 \nu_2(M,s)=\Arf(s^{-1}(p_2)).
 \label{eq:target-invariants}
\end{equation}
Then
\begin{equation}
 (\nu_3(G_3),\nu_2(G_3))=(1,0),
 \qquad
 (\nu_3(G_2),\nu_2(G_2))=(0,1),
 \label{eq:generator-invariants}
\end{equation}
because the periodic circle is nonbounding and the RR torus has odd spin
structure \cite{Atiyah1971}.  The four abstract torsion characters are
\begin{equation}
 Z_{\epsilon_3,\epsilon_2}(M,g,s)
 =(-1)^{\epsilon_3\nu_3+\epsilon_2\nu_2},
 \qquad (\epsilon_3,\epsilon_2)\in\ZZ_2^2.
 \label{eq:four-characters}
\end{equation}

\subsection{Pure heavy component on both representatives}

The common regulator can be evaluated on the part of each representative
that is already a declared unbroken gauge bundle.  On \(G_3\), take the
unbroken bundle trivial.  Its characteristic numbers are
\begin{equation}
 \int c_2(V_c)=0,\qquad \int x^2=0,\qquad \int p_1(TG_3)=0.
 \label{eq:G3-classes}
\end{equation}
On \(G_2\), take \(L=\operatorname{pr}_{S^2}^*\mathcal O(1)\), so
\(x=\operatorname{pr}_{S^2}^*u\).  Since \(u^2=0\) on \(S^2\), and the
product tangent bundle has vanishing Pontryagin number,
\begin{equation}
 \int c_2(V_c)=0,\qquad \int x^2=0,\qquad \int p_1(TG_2)=0.
 \label{eq:G2-classes}
\end{equation}
The heavy-index formula proved below therefore gives
\begin{equation}
 \boxed{\tau_h^{\rm pure}(G_3)=+1,\qquad
        \tau_h^{\rm pure}(G_2)=+1.}
 \label{eq:pure-pair}
\end{equation}
These are actual evaluations of the constant-mass unbroken-gauge plus
gravitational determinant ratio.  They do not contain a pion chiral mass or
an \(s\)-dependent fermion Yukawa family.

\subsection{No target pullback: a precise non-identifiability theorem}

Let \(\mathscr C_{\rm UV}\) be the quotient of declared smooth
\((A,\Sigma,\Phi)\) configurations by gauge transformations.  The common
quadratic regulator defines a determinant line \(\mathcal L_{\rm UV}\) over
the nondegenerate locus of \(\mathscr C_{\rm UV}\).  To ask for a target eta
pair one needs a map
\begin{equation}
 \mathfrak L:\operatorname{Map}(M_4,S^3\times S^2)
 \longrightarrow\mathscr C_{\rm UV},
 \qquad
 (g,s)\longmapsto(A,\Sigma,\Phi,\mathcal M_{\rm all\ branches}),
 \label{eq:missing-lift}
\end{equation}
and then the line \(\mathfrak L^*\mathcal L_{\rm UV}\).  The declared
microscopic action has no such map: \(g\) is a composite pion orientation,
not an elementary mass source, and \(\Phi\) does not enter the fermion mass
\(m+y\widehat\Sigma_{\bm5}\).  A Hopf connection associated with \(s\) can
specify the \(U(1)_H\) bundle used in Eq.~\eqref{eq:G2-classes}; it does not
supply the missing chiral mass on all light and heavy branches.
The continuum source in Eq.~\eqref{eq:higgsed-target-mass} does not yet
supply this map: a finite GW Dirac--Yukawa operator, its five-dimensional
kernel, and determinant orientation are still absent.  Even after those are
added, the resulting map would depend on a new Hubbard--Stratonovich source,
a no-stack convention, and a restriction to the gapped Higgs orbit, rather
than being data implicit in the original action.

\begin{theorem}[Target eta non-identifiability]
The data in Eqs.~\eqref{eq:field-list}--\eqref{eq:finite-convention} determine
the pure pair in Eq.~\eqref{eq:pure-pair}, but do not determine
\((\tau_{\rm full}(G_3),\tau_{\rm full}(G_2))\).  In particular, setting
``no torsion stack'' in the UV local scheme cannot replace the absent map
\(\mathfrak L\).
\end{theorem}
\begin{proof}
Without \(\mathfrak L\) in the original action, the pullback determinant line on which the two
holonomies would be evaluated is undefined.  This is a type obstruction,
not a difficult numerical eta sum.  Independently, suppose one arbitrary
finite lift with continuum limit Eq.~\eqref{eq:higgsed-target-mass} were
supplied.  Stacking either invertible spin character
\(I_3=(-1)^{\nu_3}\) or \(I_2=(-1)^{\nu_2}\) has zero local de Rham
curvature, changes none of the principal symbols, PV moments, gauge-bordism
character, or pure heavy index, but flips exactly one target value.  Thus
all four rows
\[
 (+,+),\quad(-,+),\quad(+,-),\quad(-,-)
\]
remain compatible with the currently declared local and pure-threshold
data.  Selection requires both \(\mathfrak L\) and an explicit microscopic
choice of the torsion invertible sector.
\end{proof}

For comparison, the circle spectral path
\(D_{S^1,\alpha_s}+t\), \(-1/4\le t\le1/4\), has one crossing for periodic
spin \(\alpha_s=0\) and none for antiperiodic spin \(\alpha_s=1/2\).  The
RR torus has one chiral zero mode modulo two, whereas the other three torus
spin structures have none.  These exact KO controls reproduce
Eq.~\eqref{eq:generator-invariants}; they realize the two abstract stacks,
not the missing \(Sp(4)\) target lift.

\begin{gate}[Full target pair]
The exact status is
\begin{equation}
 \begin{gathered}
 \texttt{pure\_heavy\_phase\_on\_G3\_computed=true},\\
 \texttt{pure\_heavy\_phase\_on\_G2\_computed=true},\\
 \texttt{conditional\_higgsed\_parity\_predicts\_plus\_plus=true},\\
 \texttt{minimal\_higgsed\_overlap\_eta\_pair\_computed=false},\\
 \texttt{original\_action\_target\_mass\_family\_lift\_defined=false},\\
 \texttt{target\_mapping\_torus\_eta\_pair\_selected=false}.
 \end{gathered}
 \label{eq:target-gates}
\end{equation}
\end{gate}

\section{Exact constant-mass heavy threshold}

\subsection{Index polynomial}

At \(m=-yV+\mu\), with \(yV>0\) and then \(\mu\downarrow0\), the branches
have signs
\begin{center}
\begin{tabular}{ccccc}
\toprule
branch & bundle & charge & limiting mass & sign\\
\midrule
light & \(V_c\otimes L\) & \(+1\) & \(\mu\) & \(+\)\\
charged heavy & \(V_c\otimes L^{-1}\) & \(-1\) & \(-2yV\) & \(-\)\\
neutral heavy & \(\bm1\) & \(0\) & \(-yV\) & \(-\)\\
\bottomrule
\end{tabular}
\end{center}
Put \(x=c_1(L)\) and use \(c_1(V_c)=0\).  The degree-four Chern characters
are
\begin{align}
 \operatorname{ch}_2(V_c\otimes L^{-1})
 &=-c_2(V_c)+x^2,
 \label{eq:charged-chern}\\
 \operatorname{ch}_2(\bm1)&=0.
 \label{eq:neutral-chern}
\end{align}
Since \(\widehat A(TM)=1-p_1(TM)/24+\cdots\), the one-flavour heavy index is
\begin{align}
 I_h
 &=\int_{M_4}\widehat A(TM_4)
   \left[\operatorname{ch}(V_c\otimes L^{-1})+1\right]_{(4)}
 \notag\\
 &=\int_{M_4}
 \left[-c_2(V_c)+x^2-\frac18p_1(TM_4)\right]\in\ZZ.
 \label{eq:heavy-index}
\end{align}
The separate gravitational coefficients are \(-p_1/12\) for the rank-two
charged branch and \(-p_1/24\) for the singlet.  Thus
Eq.~\eqref{eq:heavy-index} includes both and does not discard the neutral
fermion.

For one negative Dirac mass relative to a positive reference, paired
nonzero eigenvalues cancel and the zero modes give
\begin{equation}
 \frac{\det(\slashed D-|M|)}{\det(\slashed D+|M|)}
 =(-1)^{\dim\ker\slashed D}
 =(-1)^{\ind D^+}.
 \label{eq:mass-sign-index}
\end{equation}
The second equality holds because sum and difference of the two chiral
kernel dimensions have the same parity.  It is also the spectral-flow
evaluation of Eq.~\eqref{eq:fermion-phase}.

\begin{theorem}[Two-flavour pure threshold]
For the complete negative-mass heavy sector and \(N_f=2\),
\begin{equation}
 \tau_h^{\rm pure}(M_4,V_c,L)
 =\exp\!\left(\ii\pi N_f I_h\right)
 =\exp(2\pi\ii I_h)=1
 \label{eq:pure-heavy-theorem}
\end{equation}
on every closed spin four-manifold and every
\(SU(2)_c\times U(1)_H\) bundle.
\end{theorem}
\begin{proof}
Equation~\eqref{eq:mass-sign-index} applies additively to the charged and
neutral heavy branches, giving \(I_h\) in Eq.~\eqref{eq:heavy-index}.
The index is integral by the Atiyah--Singer theorem.  Two identical flavours
multiply it by two, so the exponential is one.  Equivalently, on a closed
spin four-manifold \(x^2\) is even and \(p_1=3\sigma\) with
\(\sigma\in16\ZZ\), while the possible odd \(c_2(V_c)\) contribution is
also doubled.
\end{proof}

The corresponding finite local threshold functional is
\begin{equation}
 S_{h,\rm loc}
 =2\pi\ii\int_{M_4}
 \left[-c_2(V_c)+x^2-\frac18p_1(TM_4)\right].
 \label{eq:local-threshold}
\end{equation}
It is an integral counterterm and exponentiates to one on closed spin
backgrounds.  On a manifold with boundary its Chern--Simons transgression
is precisely the local term paired with the APS eta section; the finite
scheme in Eq.~\eqref{eq:finite-convention} removes it.  Thus the constant
real-mass unbroken-gauge and gravitational determinant ratio is actually
matched, not merely checked for theta periodicity.

\begin{gate}[Scope of the threshold theorem]
Equation~\eqref{eq:pure-heavy-theorem} closes
\begin{equation}
 \texttt{pure\_unbroken\_gauge\_gravity\_heavy\_phase\_computed=true}.
 \label{eq:pure-threshold-true}
\end{equation}
It does not evaluate a target-dependent chiral mass family, because that
family is absent from the action.  Accordingly
\begin{equation}
 \texttt{mixed\_wzw\_heavy\_determinant\_computed=false}.
 \label{eq:mixed-threshold-false}
\end{equation}
\end{gate}

\section{Controlled all-scale mixed-WZW tests}

The mixed WZW curvature studied in the recent UV completion of
Davighi--Lohitsiri is an integral degree-five class and its coefficient is
fixed by a generalized-symmetry matching condition
\cite{DavighiLohitsiri2024}.  For the present branch weights, a uniform
exact chiral flavour action on the complete \(\bm5\) gives
\begin{equation}
 \kappa_{\rm full}=2(+1)+2(-1)+1(0)=0,
 \qquad
 \kappa_{\rm light}=+2,\qquad
 \kappa_{\rm heavy}=-2.
 \label{eq:charge-trace}
\end{equation}
The sign of a real Dirac mass does not change this charge trace: multiplying
an \(SU(2)\) chiral mass map by \(-1\) leaves
\(({-g})^{-1}\dd(-g)=g^{-1}\dd g\).

\begin{theorem}[Exact separated rank-two projector obstruction]
Let \(\mathcal U\) be an \(Sp(4)\)-invariant neighbourhood of
\(\Sigma=0\) in the adjoint representation.  Suppose
\[
 C:\mathcal U\longrightarrow\operatorname{End}_{\CC}(\bm5)
\]
is continuous and equivariant,
\(C(u\Sigma)=R_{\bm5}(u)C(\Sigma)R_{\bm5}(u)^{-1}\).
If \(C(t\Sigma_0)\) is the rank-two light-branch projector \(P_+\) for
\(t>0\), then \(C\) cannot extend to \(t=0\) while retaining a separated
rank-two branch.  More generally, a continuous equivariant chiral-mass
coefficient has
\begin{equation}
 C(0)=c\bm1_{\bm5},\qquad
 \Tr_{\bm5}\!\left(HC(0)\right)=c\Tr_{\bm5}H=0,
 \label{eq:schur-zero}
\end{equation}
whereas the Higgsed selector obeys
\begin{equation}
 \Tr_{\bm5}(HP_+)=2.
 \label{eq:projector-trace-two}
\end{equation}
\end{theorem}
\begin{proof}
The restored point is fixed by all of \(Sp(4)\).  Equivariance therefore
makes \(C(0)\) commute with the irreducible complex \(\bm5\); Schur's lemma
gives the first equation in Eq.~\eqref{eq:schur-zero}.  The eigenvalues of
\(H=h(\Sigma_0)\) are \(+1,+1,-1,-1,0\), so its trace is zero, while
Eq.~\eqref{eq:higgs-projectors} gives
\(\Tr(HP_+)=1+1=2\).  Continuity along \(t\Sigma_0\) is incompatible with
keeping this coefficient equal to two at \(t=0\).  Equivalently, if \(C\)
remains an orthogonal projector, then \(C(0)=c\bm1\) and \(C(0)^2=C(0)\)
force \(c=0\) or \(1\), of rank zero or five, never rank two.  Hence the
branch gap must close, the selector must become discontinuous, equivariance
must fail, or new UV fields must carry the missing data.
\end{proof}

This theorem is deliberately limited to an exact separated projector.  It
does not exclude continuous equivariant non-idempotent coefficients such as
\(C(\Sigma)=c\bm1+f(\lVert\Sigma\rVert)\widehat\Sigma\).  Such families can
reorganize their spectrum or close a gap at the restored point and require a
separate determinant-line analysis.

\begin{proposition}[Two controlled mass assignments]
Within the currently declared \(Sp(4)\) field/action data, the following two
assignments can be evaluated:
\begin{enumerate}[label=(\Alph*)]
\item the chiral flavour mass map acts uniformly on the full \(\bm5\), in
  which case the exact all-scale mixed coefficient is zero by
  Eq.~\eqref{eq:charge-trace}; or
\item the desired chiral flavour symmetry emerges only on the tuned light
  branch below the breaking threshold, in which case no deep-UV target
  mass-family lift exists and its APS eta pair cannot be computed.
\end{enumerate}
Neither evaluated assignment is an all-scale derivation of \(k=+2\).  No
claim is made that these two assignments exhaust all continuous equivariant
non-projector mass families.
\end{proposition}
\begin{proof}
In case (A), additivity of the family index over representation weights gives
the first equality in Eq.~\eqref{eq:charge-trace}; local four-dimensional
counterterms cannot change the integral cohomology class of the degree-five
curvature.  Restoring \(+2\) requires adding a new heavy differential
character or changing the representation/symmetry assignment.  In case
(B), the heavy Dirac masses preserve only vector flavour, so the deep-UV
operator is not a family over \(S^3\); this is exactly the type obstruction
in Eq.~\eqref{eq:missing-lift}.  New fields, a new invertible sector, or a
non-projector equivariant mass family require additional analysis outside
these two controlled assignments.
\end{proof}

The infrared orientation remains meaningful:
\begin{equation}
 n_{\rm IR}=N_cX_q=2,\qquad k_{\rm IR}=n_{\rm IR}B=+2
 \quad(X_q=+1,\ B=+1).
 \label{eq:ir-k}
\end{equation}
It is not promoted to an all-scale theorem.  Recent \(Sp\)-QCD bordism work
controls the consistency of WZW sectors but does not provide the missing
mixed target family of this particular Higgsed model \cite{Saito2024}.

\section{Deterministic verification}

The companion program
\texttt{route\_f/code/verify\_ap\_e7\_common\_sp4\_aps\_pv\_regulator.py}
performs exact or floating-point checks of:
\begin{enumerate}
\item the Euclidean Clifford algebra, gamma-five, quaternionic reference
  structure \(BB^*=-1\), and the odd transformation of the adjoint mass;
\item gamma-five Hermiticity and the two-flavour determinant square for a
  scalar-reference control;
\item all PV moments and the \(-2/\lambda^3\) asymptotic coefficient;
\item positivity and common construction of toy boson and ghost Laplace
  symbols;
\item the physical zero-cut sign, an auxiliary-cut counter-control, R/NS
  circle spectral flow, and the four abstract target characters;
\item a symbolic Chern-character expansion, spin-integrality spot checks on a
  finite characteristic-number grid, and direct evaluation of the two pure
  generator indices;
\item the conditionally derived Higgsed parity prediction, the
  \(+2-2=0\) charge trace, and explicitly labelled declared-scope gates.
\end{enumerate}
These checks verify the algebra and reproducibility.  They are not a lattice
definition of interacting \(Sp(4)\) gauge theory.
This regulator branch is background-field/topological and does not import or
depend on the relaxed \(B=1\) soliton profile or its checksum.

\section{Final gate ledger and ordered continuation}

\begin{longtable}{@{}L{0.36\textwidth}L{0.16\textwidth}L{0.39\textwidth}@{}}
\toprule
claim & status & exact boundary\\
\midrule
\endfirsthead
\toprule
claim & status & exact boundary\\
\midrule
\endhead
Euclidean gamma and reference-reality data & \status{closed with scope} &
explicit Clifford representation, \(J^2=-1\), scalar-reference Kramers
sector; fixed adjoint mass transforms with \(\Sigma\mapsto-\Sigma\)\\
Full common quadratic gauge--scalar--ghost--fermion PV amplitude &
\status{closed} & exact gauge complex, all mixing retained, shared four-node
superdeterminant and finite local scheme\\
Five-dimensional APS problem & \status{closed for UV field paths} & zero-cut
physical phase with gapped endpoints; arbitrary cut retained only for the
Fredholm index; complementary APS domains and self-adjoint double\\
Nonperturbative continuum \(Sp(4)\) measure & \status{open} & no constructive
interacting cutoff removal, global gauge slice, or reflection-positivity proof\\
Pure unbroken gauge plus gravitational heavy threshold & \status{closed} &
two-flavour phase \(\exp(2\pi\ii I_h)=1\), local integral counterterm explicit\\
Pure heavy phase on \(G_3,G_2\) & \status{closed} & \((+1,+1)\) for the
constant-mass unbroken bundle component\\
Conditional Higgsed GW source & \status{source only} & \((+1,+1)\) is a
copy-counting prediction; finite Dirac--Yukawa operator, five-dimensional
target kernel, determinant orientation, and eta computation remain open\\
Full target eta pair & \status{undefined} & no all-branch target mass map or
pullback determinant line in the original action; two torsion stacks prove
non-identifiability\\
Uniform exact-flavour assignment & \status{evaluated, gives zero} &
complete-irrep coefficient \(+2-2=0\); not an exhaustive classification of
equivariant non-projector families\\
Light emergent-flavour \(k_{\rm IR}=+2\) & \status{derived only in IR} & no
deep-UV target-family APS matching\\
Degree-one Route-E portal & \status{forbidden} & neither target matching nor
full \(Sp(4)\) dynamical lane is closed\\
\bottomrule
\end{longtable}

The machine-readable final flags are
\begin{equation}
 \begin{gathered}
 \texttt{euclidean\_gamma\_reference\_reality\_data\_specified=true},\\
 \texttt{exact\_separated\_rank\_two\_projector\_obstruction\_proven=true},\\
 \texttt{common\_quadratic\_aps\_pv\_regulator\_specified=true},\\
 \texttt{five\_dimensional\_aps\_domain\_specified=true},\\
 \texttt{physical\_fermion\_phase\_zero\_cut\_fixed=true},\\
 \texttt{fixed\_adjoint\_mass\_kramers\_reality\_proven=false},\\
 \texttt{full\_nonperturbative\_sp4\_measure\_constructed=false},\\
 \texttt{pure\_unbroken\_gauge\_gravity\_heavy\_phase\_computed=true},\\
 \texttt{conditional\_higgsed\_gw\_source\_specified=true},\\
 \texttt{finite\_higgsed\_overlap\_dirac\_yukawa\_operator\_defined=false},\\
 \texttt{five\_dimensional\_higgsed\_target\_kernel\_defined=false},\\
 \texttt{minimal\_higgsed\_overlap\_eta\_pair\_computed=false},\\
 \texttt{conditional\_higgsed\_parity\_predicts\_plus\_plus=true},\\
 \texttt{original\_action\_target\_mass\_family\_lift\_defined=false},\\
 \texttt{target\_mapping\_torus\_eta\_pair\_selected=false},\\
 \texttt{mixed\_wzw\_heavy\_determinant\_computed=false},\\
 \texttt{k\_plus\_two\_all\_scale\_matched=false},\\
 \texttt{sp4\_aps\_pv\_lane\_closed=false},\\
 \texttt{degree\_one\_route\_e\_portal\_started=false},\\
 \texttt{physics\_promotion\_authorized=false}.
 \end{gathered}
 \label{eq:final-flags}
\end{equation}

The next discriminating calculations are:
\begin{enumerate}
\item write the exact finite GW Dirac--Yukawa operator, its determinant
  orientation, and its five-dimensional domain-wall target kernel before
  promoting the conditional \((+1,+1)\) parity count to an eta result;
\item derive the Hubbard--Stratonovich source in
  Eq.~\eqref{eq:higgsed-target-mass} from the original interacting action,
  or replace it by new UV fields that evade the exact-projector obstruction;
\item if exact uniform flavour is retained, add a specified heavy
  differential character of coefficient \(+2\) and calculate both of its
  torsion signs rather than choosing them;
\item alternatively, formulate a symmetry-emergence matching theorem across
  the \(Sp(4)\) breaking scale, including the operator that creates the
  light-only chiral determinant line;
\item only after one of these closes, combine it with the independent
  relaxed-soliton or \(SO(3)\)/FR determinant-line routes.  The portal remains
  last.
\end{enumerate}

\bibliographystyle{unsrt}
\bibliography{ap_e7_common_sp4_aps_pv_regulator}

\end{document}
